% Motivation and learning outcomes.
% The outcomes are the outcomes of Day 12 in the syllabus.
\begin{frame}{Why this day matters}
  \begin{itemize}
    \item On Day 11 the boards sent video. A system that sends only results
          needs less bandwidth and keeps the images on the device. Today the
          boards send small messages: a class, a reading, a command.
    \item Some decisions cannot wait for a server: a door that must close, a
          motor that must stop. The device must decide alone, also when the
          network or the internet fails.
    \item MQTT is the common protocol for such messages. A broker receives
          each message and gives it to the devices that asked for its topic.
          The protocol has settings for lost messages, lost devices, and the
          last value.
    \item In the lab you install a broker on the Raspberry Pi, send IMU
          readings and model results from the XIAO, apply a local rule that
          sends a command back, and show that the system works with no
          internet and loses no message.
  \end{itemize}
\end{frame}

\begin{frame}{Learning outcomes}
  After this day you can:
  \begin{itemize}
    \item Design a cascade that makes decisions locally.
    \item Explain MQTT topics, quality of service levels, retained messages,
          and the last will message.
    \item Publish inference results from a microcontroller.
    \item Build a system that continues to work without an internet
          connection.
  \end{itemize}
\end{frame}
